% Atur numbering halaman ke Roman dan mulai dari i
\pagenumbering{roman}
\setcounter{page}{1}

% Atur spacing untuk konten lembar pengesahan
\setlength{\parskip}{8pt}

% Macro untuk baris examiner/advisor (mengurangi redundansi)
\newcommand{\examrow}[3]{#1 & #2 \\ NIP: #3 & \\ & ................................... \\ & \\ & \\}

% Macro untuk baris dengan spacing (menghilangkan vspace repetitif)
\newcommand{\spline}[1]{\vspace{6pt} \\ #1}

\begin{center}
  \fontsize{14}{16.8}\selectfont
  \textbf{LEMBAR PENGESAHAN}
\end{center}
\vspace{10pt}

\begin{center}
  \textbf{\tatitle{}}
\end{center}

\begingroup
% Pemilihan font ukuran small
\small

\begin{center}
  \vspace{6pt}
  \textbf{TUGAS AKHIR} \\
  \vspace{6pt}
  Diajukan untuk memenuhi salah satu syarat \spline{memperoleh gelar Sarjana Komputer pada} \spline{Program Studi S-1 \studyprogram{}} \spline{Departemen \department{}} \spline{Fakultas \faculty{}} \spline{Institut Teknologi Sepuluh Nopember}
\end{center}

\begin{center}
  Oleh: \textbf{\name{}} \\[6pt]
   NRP. \nrp{}
\end{center}

\begin{center}
  Disetujui oleh Tim Penguji Tugas Akhir: 
\end{center}

\begingroup
% Menghilangkan padding
\setlength{\tabcolsep}{0pt}

\noindent
\begin{tabularx}{\textwidth}{X l}
  \examrow{\advisor{}}{\text{(Pembimbing I)}}{\advisornip{}}
  \examrow{\coadvisor{}}{\text{(Pembimbing II)}}{\coadvisornip{}}
  \examrow{\examinerone{}.}{\text{(Penguji I)}}{\examineronenip{}}
  \examrow{\examinertwo{}.}{\text{(Penguji II)}}{\examinertwonip{}}
  % \examrow{\examinerthree{}.}{\text{(Penguji III)}}{\examinerthreenip{}}
\end{tabularx}
\endgroup

\begin{center}
  \textbf{\MakeUppercase{\place{}}} \\
  \textbf{\MONTH{}}, \textbf{\the\year{}}
\end{center}
\endgroup

